% ch05_u03 —— 原书 p.210–p.219（5.3.18–5.3.31、5.4 节首与 5.4.1–5.4.7）
\SourcePage{210}

\bdp{5.3.18} 对 $M\in\mathcal{P}$，置 $h^{\cdot}(M):=H^{\cdot}_{DR}(\mathcal{GR},M)$。这是 $\mathbb{Z}$-分次向量空间；记 $h^{\varepsilon}(M)$ 为相应的 $\mathbb{Z}/2\mathbb{Z}$-分次向量空间。

考虑投影 $p:\mathcal{GR}_X\to X$。$X$ 上的 $\mathcal{D}$-模 $H^ap_*(M_X)$ 是常值的，即同构于 $\omega_X$ 的若干拷贝之和（回忆：我们处理的是右 $\mathcal{D}$-模）。相应的纤维是 $h^{\cdot}(M)$：对任意 $x\in X$，有 $H^{\cdot}i_x^{!}p_*(M_X)=h^{\cdot}(M)$，其中 $i_x$ 是嵌入 $\{x\}\hookrightarrow X$。

\bdp{5.3.19} \emph{命题.}（[MV]）对 $\mathcal{P}$ 中对象的任意族 $\{M_j\}_{j\in J}$，$X^J$ 上的 $\mathcal{D}$-模 $H^ap_*^{(J)}(\boxtimes M_{jX})$ 是常值的。\hfill$\Box$

对任意 $(x_j)\in X^J$，有
\begin{equation}\tag{270}
H^{\cdot}i^{!}_{(x_j)}p_*^{(J)}(\boxtimes M_{jX})=\otimes h^{\cdot}(M_j).
\end{equation}

由 5.3.18，这对 $(x_j)\in U^{(J)}$ 显然；然后用 5.3.19。
% [Tutor Note: 原书 p.210 此处作「then use 5.3.19」，与正在证明的命题同号，底本如此，保留原样。]

\bdp{5.3.20} 对 $(x_j)\in X\subset X^J$，（270）给出典范同构 $h^{\cdot}(\circledast M_j)=\otimes h^{\cdot}(M_j)$，它显然与「约束」（266）相容。于是
\begin{equation}\tag{271}
h^{\cdot}:\mathcal{P}\to\operatorname{Vect}^{\cdot},\qquad h^{\varepsilon}:\mathcal{P}\to\operatorname{Vect}^{\varepsilon}
\end{equation}
是\emph{张量函子}。这里 $\operatorname{Vect}^{\cdot}$ 是带「超」交换约束的 $\mathbb{Z}$-分次向量空间的张量范畴，$\operatorname{Vect}^{\varepsilon}$ 是 $\mathbb{Z}/2\mathbb{Z}$-分次向量空间的类似张量范畴。

\bdp{5.3.21} 可以扭 $\mathcal{P}$ 上的张量结构，以摆脱超向量空间。为此注意：$\mathcal{P}$ 的对象带典范的 $\mathbb{Z}/2\mathbb{Z}$-分次 $\varepsilon$，它由支撑各分支的奇偶性给出（见 4.5.10）。该分次与 $\boxtimes$ 相容。

记 $\mathcal{P}^{\natural}$ 为 $\mathcal{P}^{\varepsilon}:=\mathcal{P}\otimes\operatorname{Vect}^{\varepsilon}$ 中偶对象的全子范畴（按 $\mathbb{Z}/2\mathbb{Z}$-分次的张量积取）。它是 $\mathcal{P}^{\varepsilon}$ 的张量子范畴。「遗忘分次」函子 $o_{\varepsilon}:\operatorname{Vect}^{\varepsilon}\to\operatorname{Vect}$ 给出等价 $\mathcal{P}^{\natural}\xrightarrow{\sim}\mathcal{P}$。这是幺半范畴的等价（即它与张量积、结合约束相容）；
\SourcePage{211}
$\mathcal{P}$ 与 $\mathcal{P}^{\natural}$ 的交换约束 $A\otimes B\cong B\otimes A$ 相差 $(-1)^{p(A)p(B)}$。
% [Tutor Note: 原书 p.211 起数处出现作者彼此提醒用的大写插语（WHAT ABOUT GINZBURG ??、WHAT ABOUT THE SIGN??? 等），译文按其语气译出，底本如此。]

由 5.3.3(ii)，函子 $h^{\varepsilon}$ 与 $\mathbb{Z}/2\mathbb{Z}$-分次相容。因此它定义张量函子
\begin{equation}\tag{272}
h:\mathcal{P}^{\natural}\to\operatorname{Vect}.
\end{equation}
注意 $h$ 带典范的 $\mathbb{Z}$-分次，因记号滥用，我们仍把它记作 $h^{\cdot}$。于是 $h^{\cdot}$ 是 $\mathcal{P}^{\natural}$ 上的张量函子，取值于带普通（非超）交换约束的分次向量空间张量范畴。

\bdp{5.3.22} 据 [MV]（GINZBURG 怎么办？？），张量范畴 $\mathcal{P}^{\natural}$ 是刚性的，即每个对象都有 [Del91] §2.1.2 意义下的对偶（对偶对象在 [MV] 中有显式构造）。张量函子（272）是 $\mathbb{C}$-线性的且正合\footnote{正合性显然，因为 $\mathcal{P}^{\natural}$ 半单。Mirković 与 Vilonen [MV] 必须证明正合性，因为他们希望其证明对任意交换环上的反常层都有效。}，故它是 [Del91] 意义下的纤维函子。于是由一般的 Tannakian 形式化，（272）诱导张量范畴 $\mathcal{P}^{\natural}$ 与 $\operatorname{Rep}(\operatorname{Aut}^{\otimes}h)$ 之间的等价；这里 $\operatorname{Aut}^{\otimes}h$ 记 $h$ 的张量自同构群概形，$\operatorname{Rep}$ 指有限维表示范畴。据 [MV]，有同构 $\varkappa:{}^{L}G\xrightarrow{\sim}\operatorname{Aut}^{\otimes}h$，故可把上述等价改写为
\begin{equation}\tag{273}
h:\mathcal{P}^{\natural}\xrightarrow{\sim}\operatorname{Rep}{}^{L}G.
\end{equation}
这里 ${}^{L}G$ 是朗兰兹对偶群：它是半单群，带固定的嘉当环面 ${}^{L}H\subset{}^{L}G$、相应的根数据与 $G$ 的根数据之对偶的等同，以及对每个单负根 $\alpha$ 固定的一组非零向量 $y_{\alpha}\in({}^{L}\mathfrak{g})^{\alpha}$。

\SourcePage{212}

\bdp{5.3.23} 我们将通过列举 ${}^{L}G$ 在 $h$ 上作用的若干性质，来定义一个\emph{典范同构}
\begin{equation}\tag{274}
\varkappa:{}^{L}G\xrightarrow{\sim}\operatorname{Aut}^{\otimes}h,
\end{equation}
这些性质唯一确定 $\varkappa$。

(i) 记
\begin{equation}\tag{275}
t\mapsto t^{2\rho}
\end{equation}
为相应于 $G$ 的权 $2\rho$ 的态射 $\mathbb{G}_m\to{}^{L}H$。则 $t^{2\rho}$ 以乘以 $t^{-a}$ 的方式作用在 $h^a$ 上（故单参数子群（275）的作用对应于 $h$ 的分次 $h^{\cdot}$）。由 (i) 推出，${}^{L}H$ 在 $h$ 上的作用保持 $h$ 的分次。

(ii) 对任意 $\chi\in P_+({}^{L}G)$，群 ${}^{L}H$ 通过特征标 $\chi$ 作用在 $h^{\min}(I_{\chi})=h^{-\dim\operatorname{Orb}_{\chi}}(I_{\chi})$ 上。这意味着不可约 ${}^{L}G$-模 $h(I_{\chi})$ 的最高权等于 $\chi$。

\emph{注.} 由于 $\dim\operatorname{Orb}_{\chi}=\langle\chi,2\rho\rangle$，(i) 与 (ii) 之间没有矛盾。

性质 (i) 与 (ii) 可在 [MV] 中找到。它们唯一确定（274）在 ${}^{L}H$ 上的限制。故 (274) 由 (i)、(ii) 唯一确定到 ${}^{L}H$-共轭。我们用下述性质规范化（274）。

(iii) 设 $c\in(\operatorname{Sym}^2\mathfrak{g}^*)^G$ 是 $\mathfrak{g}$ 上的不变双线性型（约化情形下是 $[\mathfrak{g},\mathfrak{g}]$ 上的呢？？？）。置
\begin{equation}\tag{276}
f_c:=\frac{???}{2}\sum_{\alpha}c(\alpha,\alpha)y_{\alpha}\in{}^{L}\mathfrak{g}
\end{equation}
% [Tutor Note: 原书 p.212 公式（276）的分子印作「???」，系底本缺失的记号，底本如此，保留原样。]
（表达式 $c(\alpha,\alpha)$ 有意义，因为 $\alpha\in({}^{L}\mathfrak{h})^*=\mathfrak{h}\subset\mathfrak{g}$。）则李代数元素 $f_c$ 以乘以 $\nu(c)$ 的方式作用在 $h(M)=H^{\cdot}_{DR}(\mathcal{GR},M)$ 上，其中 $M\in\mathcal{P}^{\natural}=\mathcal{P}$，而
\begin{equation}\tag{277}
\nu:(\operatorname{Sym}^2\mathfrak{g}^*)^G\to H^2_{DR}(\mathcal{GR})
\end{equation}

\SourcePage{213}

式（277）是标准态射，其定义将在 5.3.24 中回顾。

\emph{注.} (iii) 由 V. Ginzburg [Gi95] 以略有不同的形式给出。事实上，他用类似方式描述了 $f_c$ 在 ${}^{L}\mathfrak{g}$ 中的整个中心化子在 $h$ 上的作用。

\bdp{5.3.24} 本小节（读者可跳过）定义典范态射（277）。我们采用如下的 ad hoc 定义：对任意归纳概形 $Z$，置 $H^a_{DR}(Z):=\varprojlim H^a(Y,\Omega_Y)$，其中 $Y$ 跑遍 $Z$ 的所有闭子概形，$\Omega_Y$ 是 $Y$ 的 de Rham 复形（最朴素意义下）。为定义 $\nu$，姑且设（简简单单两次？？糟糕的文风）$G$ 半单\footnote{任意 $G$ 的 $\nu$ 之定义留作习题。}。则投影 $G(K)\to\mathcal{GR}$ 诱导同构 $H^2_{DR}(\mathcal{GR})\xrightarrow{\sim}H^2_{DR}(G(K))$（事实上，该投影是 $G(\mathcal{O})$-torsor，$G(\mathcal{O})$ 连通，且 $H^1_{DR}(G(\mathcal{O}))=H^2_{DR}(G(\mathcal{O}))=0$）。现在 $c$ 定义了 $\mathfrak{g}\otimes K$ 上的 Kac–Moody 余圈 $u,v\mapsto\operatorname{Res}_{t=0}c(du,v)$。设 $\omega_c$ 是 $G(K)$ 上相应的右不变闭 2-形式。其类经上述同构之逆映射所得的像为 $\nu(c)\in H^2_{DR}(\mathcal{GR})$。符号怎么办？？？

\emph{注.} 在 5.3.23(iii) 中，我们用了 $H^{\cdot}_{DR}(\mathcal{GR})$ 在 $H^{\cdot}_{DR}(\mathcal{GR},M)$ 上的作用，其中 $M$ 是 $\mathcal{GR}$ 上的 $\mathcal{D}$-模。该作用定义如下。考虑 $\Omega^{\cdot}$-复形 $\Omega M$（见 7.11.13）。则 $H^{\cdot}_{DR}(\mathcal{GR},M)=\varinjlim H^{\cdot}(Y,\Omega M_{(Y)})$，其中 $Y$ 跑遍 $\mathcal{GR}$ 的所有子概形。现在 $\Omega M_{(Y)}$ 是 $Y$ 上的 $\Omega$-复形，故 $H^{\cdot}(Y,\Omega_Y)$ 作用在 $H^{\cdot}(Y,\Omega M_{(Y)})$ 上。因此 $H^{\cdot}_{DR}(\mathcal{GR})$ 作用在 $H^{\cdot}_{DR}(\mathcal{GR},M)$ 上。

\bdp{5.3.25} 5.3.23 中给出的典范同构（274）的简要刻画已足够我们使用。想更好理解（274）的读者可读 ???-??? 与 [MV]。
% [Tutor Note: 原书 p.213 此处作「the reader may read ???-???」，交叉引用印作问号，底本如此，保留原样。]

\bdp{5.3.26}

\SourcePage{214}

\emph{注.} 回忆（见 4.5.9）$\mathcal{GR}$ 的连通分支由 $Z({}^{L}G)^{\vee}$ 的元素标号；这里 $Z({}^{L}G)^{\vee}$ 是中心 $Z({}^{L}G)\subset{}^{L}G$ 的特征标群。$\mathcal{GR}$ 中相应于 $\zeta\in Z({}^{L}G)^{\vee}$ 的连通分支记作 $\mathcal{GR}_{\zeta}$。支撑分解 $D(\mathcal{GR})=\prod D(\mathcal{GR}_{\zeta})$、$\mathcal{P}=\oplus\mathcal{P}_{\zeta}$ 定义了 $h$ 上的 $Z({}^{L}G)^{\vee}$-分次，即 $Z({}^{L}G)$-作用。该作用与 ${}^{L}G$-作用所诱导者一致。

本节余下部分解释上述构造如何与过到 ${}^{L}G$ 的 Levi 子群相容。当该子群为 ${}^{L}H\subset{}^{L}G$ 时，这归结为 Mirković–Vilonen 给出的 ${}^{L}H$ 在纤维函子 $h$ 上作用的显式描述。

\bdp{5.3.27} 设 $P\subset G$ 是抛物子群，$N_P\subset P$ 是其幂幺根，$F:=P/N_P$ 是 Levi 群。$F$ 与 $G$ 的嘉当环面以显然方式等同，$F$ 的根数据是 $G$ 的根数据的子集。于是对标准环面 ${}^{L}H\subset{}^{L}F\subset{}^{L}G$，${}^{L}F$ 是 ${}^{L}G$ 的 Levi 子群。从而 $Z({}^{L}G)\subset Z({}^{L}F)$。

我们将定义典范张量函子
\begin{equation}\tag{278}
r_P^{\natural}:\mathcal{P}_G^{\natural}\to\mathcal{P}_F^{\natural},
\end{equation}
它经由等价 $h_G$、$h_F$ 对应于显然的限制函子 $r^{GF}:\operatorname{Rep}{}^{L}G\to\operatorname{Rep}{}^{L}F$。

\bdp{5.3.28} 图表 $G\leftarrow P\rightarrow F$ 给出相应仿射 Grassmann 流形的态射
\begin{equation}\tag{279}
\mathcal{GR}^G\xleftarrow{i}\mathcal{GR}^P\xrightarrow{\pi}\mathcal{GR}^F.
\end{equation}
这里 $\pi$ 是形式光滑的归纳仿射满投影。其纤维是 $N_P(K)$-轨道。故 $\pi$ 给出 $\mathcal{GR}^P$ 与 $\mathcal{GR}^F$ 的连通分支集之间的双射。对任意 $\zeta\in Z({}^{L}F)^{\vee}$，记 $\mathcal{GR}^P_{\zeta}$ 为相应的分支。则 $i$ 的限制 $i_{\zeta}:\mathcal{GR}^P_{\zeta}\hookrightarrow\mathcal{GR}^G$ 是局部闭嵌入；其像落在 $\mathcal{GR}^G_{\bar{\zeta}}$ 中，其中 $\bar{\zeta}:=\zeta|_{Z({}^{L}G)}$。归纳概形 $\mathcal{GR}^P_{\zeta}$ 构成 $\mathcal{GR}^G$ 的分层（即对任意闭
\SourcePage{215}
子概形 $Y\subset\mathcal{GR}^G$，交 $Y_{\zeta}:=Y\cap\mathcal{GR}^P_{\zeta}$ 构成 $Y$ 的分层）。

置 $\rho_{GF}:=\rho_G-\rho_F\in\mathfrak{h}^*$。由于 $2\rho_{GF}$ 是 $F$ 的特征标（伴随作用在 $\mathfrak{n}_P$ 上的行列式），可把它看作 $Z({}^{L}F)\subset{}^{L}H$ 的单参数子群。于是对如上每个 $\zeta$，有一个整数 $\langle\zeta,2\rho_{GF}\rangle$。记 $\mathcal{GR}^F_n$ 为满足 $\langle\zeta,2\rho_{GF}\rangle=n$ 的分支 $\mathcal{GR}^F_{\zeta}$ 之并。有相应的分解 $D(\mathcal{GR}^F)=\prod D(\mathcal{GR}^F_n)$、$\mathcal{P}^F=\oplus\mathcal{P}^F_n$。置 $\mathcal{P}^{F'}=\oplus\mathcal{P}^F_n[-n]\subset D(\mathcal{GR}^F)$。如同 5.3.18，对 $M\in\mathcal{P}^{F'}$ 置 $h_F(M)=H^{\cdot}(\mathcal{GR}^F,M)\in\operatorname{Vect}^{\cdot}$。

\bdp{5.3.29} \emph{命题.}
\begin{enumerate}[label=(\roman*),leftmargin=2.2em]
\item 函子 $r_D^{GF}:=\pi_*i^{!}:D(\mathcal{GR}^G)\to D(\mathcal{GR}^F)$ 把 $\mathcal{P}^G$ 送入 $\mathcal{P}^{F'}$，故有
\begin{equation}\tag{280}
r_P^{GF}:\mathcal{P}^G\to\mathcal{P}^{F'}.
\end{equation}
\item 有函子的典范等同
\begin{equation}\tag{281}
h_G^{\cdot}=h_F^{\cdot}r_P^{GF}:\mathcal{P}^G\to\operatorname{Vect}^{\cdot}.
\end{equation}
\end{enumerate}

\emph{证明.} 先设 $P=B$ 是 Borel 子群。则 $F=H$，且 $\mathcal{GR}^H_{red}=({}^{L}H)^{\vee}$，故 $\mathcal{GR}^H$ 上的 $\mathcal{D}$-模等同于 $({}^{L}H)^{\vee}$-分次向量空间，即 ${}^{L}H$-模。分层 $\mathcal{GR}^B_{\zeta}$ 正是 $\mathcal{GR}^G$ 上的 $N_B(K)$-轨道。因此 5.3.29 就是 [MV] 的关键定理。

回忆等同（281）的构造如下（见 [MV]）。设 $\overline{\mathcal{GR}}^B_n\subset\mathcal{GR}^G$ 为 $\mathcal{GR}^B_n:=\pi^{-1}(\mathcal{GR}^H_n)$ 在 $\mathcal{GR}^G$ 中的闭包。则各 $\overline{\mathcal{GR}}^B_n$ 给出 $\mathcal{GR}^G$ 上的递减滤链。对任意 $M\in\mathcal{P}^G$，显然的态射
\[
h^n_Hr_P^{GH}(M)=H^n(\mathcal{GR}^B_n,i^{!}M)\leftarrow H^n_{\overline{\mathcal{GR}}^B_n}(\mathcal{GR}^G,M)\rightarrow H^n(\mathcal{GR}^G,M)=h^n_G(M)
\]
都是同构。它们的复合是（281）。

现在设 $P$ 是任意抛物子群。取 Borel 子群 $B\subset P$，于是 $B_F:=B/N_P\cap B$ 是 $F$ 的 Borel 子群。考虑函子 $r_D^{GH}:D(\mathcal{GR}^G)\to D(\mathcal{GR}^H)$ 与 $r_D^{FH}:D(\mathcal{GR}^F)\to D(\mathcal{GR}^H)$。由基变换，有函子的典范等同 $r_D^{GH}=r_D^{FH}r_D^{GF}$。设
\SourcePage{216}
$\mathcal{P}^{H'}\subset D(\mathcal{GR}^H)$ 是由 $B\subset G$ 所定义的范畴，故已知 $r_D^{GH}(\mathcal{P}^G)\subset\mathcal{P}^{H'}$，且（因 $\rho_{GF}=\rho_G-\rho_F$）有 $r_D^{FH}(\mathcal{P}^{F'})\subset\mathcal{P}^{H'}$。

函子 $r_P^{GH}:\mathcal{P}^{F'}\to\mathcal{P}^{H'}$ 忠实（因为差一个平移，它与 $h_F^{\cdot}$ 一致）。于是使所有 $H^iT$ 都属于 $\mathcal{P}^F$ 的对象 $T\in D(\mathcal{GR}^F)$ 属于 $\mathcal{P}^{F'}$，当且仅当 $r_D^{FH}(T)\subset\mathcal{P}^{H'}$。
% [Tutor Note: 原书 p.216 此处作「since up to shift if coincides with h^·_F」，其中「if」疑为「it」之误，底本如此，保留原样。]

把该注记用于 $T=r_D^{GF}(M)$（$M\in\mathcal{P}^G$），可见 $r_D^{GF}(M)\in\mathcal{P}^{F'}$，这正是 5.3.29(i)。我们还知道 $h_G^{\cdot}(M)=h_H^{\cdot}(r_P^{GH}(M))=h_H^{\cdot}(r_P^{FH}(M))=h_F^{\cdot}r^{GF}(M)$，这就是等同 5.3.29(ii)。至于它不依赖辅助选取的 Borel 子群 $B\subset P$，留作习题。\hfill$\Box$

\bdp{5.3.30} 范畴 $\mathcal{P}^{F'}$ 带典范张量结构（由用于 $\mathcal{P}^F$ 的同样构造定义）。函子 $r_P^{GF}:\mathcal{P}^G\to\mathcal{P}^{F'}$ 以典范方式为张量函子。事实上，（279）是手征半群的态射，故可考虑相应的函子 $r_D^{GF}:=\pi_*i^{!}:D(\mathcal{GR}^G_{X^J})\to D(\mathcal{GR}^F_{X^J})$。留作习题（提示：用 5.3.19）：对 $M_j\in\mathcal{P}^G$，该函子把 $\boxtimes M_j$ 送到 $\boxtimes r_D^{GF}(M_j)$（记号见 5.3.14）。由于它（由基变换）还与函子 $\widetilde{\Delta}^{(J)!}$ 交换，我们得到所需的张量积相容性。如同 5.3.19 可见，（281）是张量函子的同构。

最后，如同 5.3.21，把张量范畴 $\mathcal{P}^G$ 换成 $\mathcal{P}^{G\natural}$。由于 $\rho_{GF}=\rho_G-\rho_F$，可见 $r_P^{GF}$ 给出张量函子 $r^{GF}:\mathcal{P}^{G\natural}\to\mathcal{P}^{F'\natural}$，它与纤维函子 $h_G$、$h_F$ 相容。它定义了态射 $r:\operatorname{Aut}^{\otimes}h_F\to\operatorname{Aut}^{\otimes}h_G$。

\bdp{5.3.31} \emph{引理.} 态射 $\varkappa_G^{-1}r\varkappa_F:{}^{L}F\to{}^{L}G$ 与 5.3.27 中的典范嵌入一致。\hfill$\Box$

\bdsec{5.4. 主定理 II：从局部到整体}

本节给出 5.2 中主定理的精确版本，并证明局部主定理蕴含整体主定理。我们本质地用到第 7 章的「Hecke 模式」。要理解发生了什么，必须（且几乎充分）阅读 7.1.1 与 7.9.1。

\SourcePage{217}

\bdp{5.4.1} 我们从 Hecke 特征 $\mathcal{D}$-模的定义开始。考虑带 $\operatorname{Aut}\mathcal{O}$ 作用的对 $(G(K),G(\mathcal{O}))$。设 $\mathcal{H}$ 是相应的 $(\operatorname{Der}\mathcal{O},\operatorname{Aut}^0\mathcal{O})$-等变 Hecke 范畴，定义见 7.9.2\footnote{我们的 $(G(K),G(\mathcal{O}))$、$(\operatorname{Der}\mathcal{O},\operatorname{Aut}^0\mathcal{O})$ 分别是 7.9.2 的 $(G,K)$、$(\mathfrak{l},P)$。}。由于 $\mathcal{P}$ 的每个对象都以典范方式为 $\operatorname{Aut}\mathcal{O}$-等变 $\mathcal{D}$-模\footnote{据 5.3.4，$\mathcal{P}$ 的每个对象带唯一的强 $\operatorname{Aut}^0\mathcal{O}$-作用，它与强 $\operatorname{Aut}\mathcal{O}$-作用相同。}，$\mathcal{P}$ 是 $\mathcal{H}$ 的全子范畴。由定义可知，嵌入 $\mathcal{P}\to\mathcal{H}$ 是幺半函子。

考虑 $X$ 上典范的 $\operatorname{Aut}\mathcal{O}$-结构 $X^{\wedge}$（见 2.6.5）与 2.8.3 中定义的 $X^{\wedge}$ 上的概形 $M^{\wedge}$；它带 $\operatorname{Aut}\mathcal{O}\times G(K)$ 的典范作用（见 2.8.3–2.8.4）。商叠 $(\operatorname{Aut}^0\mathcal{O}\times G(\mathcal{O}))\backslash M^{\wedge}$ 等于 $\operatorname{Bun}_G\times X$。于是我们来到 7.9.1、7.9.4\footnote{我们的 $X^{\wedge}$ 与 $M^{\wedge}$ 是 7.9.4 的 $X^{\wedge}$ 与 $Y^{\wedge}$。} 的设定。故 $\mathcal{H}$ 作用在 $D(\operatorname{Bun}_G\times X)$ 上。因此 $D(\operatorname{Bun}_G\times X)$ 是 $\mathcal{P}$-模。借助 Satake 等价（273）把幺半范畴 $\mathcal{P}$\footnote{本节（5.4.6 的注除外）只用 $\mathcal{P}$ 的幺半结构（交换约束不起作用）。故可把 $\mathcal{P}$ 与 $\mathcal{P}^{\natural}$ 等同。} 与 $\operatorname{Rep}{}^{L}G$ 等同，便得到 $\operatorname{Rep}{}^{L}G$ 在 $D(\operatorname{Bun}_G\times X)$ 上的典范作用，称为 Hecke 作用。记作 $\circledast$。

注意 $D(\operatorname{Bun}_G\times X)$ 还带张量范畴 $\operatorname{Vect}^{\nabla}(X)$ 的显然作用；$\operatorname{Vect}^{\nabla}(X)$ 由 $X$ 上带联络的向量丛构成（事实上，由 $X$ 上无挠左 $\mathcal{D}$-模构成的更大张量范畴亦可）。我们把它记作 $\otimes$。它与 Hecke 作用交换，故 $D(\operatorname{Bun}_G\times X)$ 是 $(\operatorname{Rep}{}^{L}G,\operatorname{Vect}^{\nabla}(X))$-双模。

设 $\mathfrak{F}$ 是 $X$ 上带联络的 ${}^{L}G$-丛。它给出张量函子 $\operatorname{Rep}{}^{L}G\to\operatorname{Vect}^{\nabla}(X)$，$V\mapsto V_{\mathfrak{F}}$，从而给出 $\operatorname{Rep}{}^{L}G$ 在 $D(\operatorname{Bun}_G\times X)$ 上相应的作用。

\bdp{5.4.2} 设 $M$ 是 $\operatorname{Bun}_G$ 上的 $\mathcal{D}$-模。设 $M_{(X)}\in\mathcal{M}(\operatorname{Bun}_G\times X)$ 是 $M$ 的拉回。假设对任意 $V\in\operatorname{Rep}{}^{L}G$，给定自然的
\SourcePage{218}
同构 $\alpha_V:V\circledast M_{(X)}\xrightarrow{\sim}M_{(X)}\otimes V_{\mathfrak{F}}$（特别地，$V\circledast M_{(X)}$ 是 $\mathcal{D}$-模，而不仅是导出范畴的对象）。若对任意 $V_1,V_2\in\operatorname{Rep}{}^{L}G$ 有 $\alpha_{V_1\otimes V_2}=\alpha_{V_1}\circ(V_1\circledast\alpha_{V_2})$，则称这些 $\alpha_V$ 在 $M$ 上定义了\emph{Hecke $\mathcal{F}$-特征模结构}。称这样的 $(M,\alpha_V)$（或简记 $M$）为\emph{Hecke $\mathfrak{F}$-特征模}。
% [Tutor Note: 原书 p.218 前两处 oper 记号印作花体 ℱ，同页后文与 5.4.3 中均印作 𝔉，底本混用，保留原样。]

\emph{注.} 对 $X$ 上任意 ${}^{L}G$-局部系统 $\mathcal{F}$，我们希望定义 Hecke $\mathfrak{F}$-特征模的三角范畴\footnote{当然，在 Hecke 特征模的上述定义中，$M$ 可取 $D(\operatorname{Bun}_G)$ 的任意对象，而不只是 $\mathcal{D}$-模。但在此一般性下该定义看来并不合理（这类对象不构成三角范畴）。}。

下述定理是定理 5.2.6 的精确版本。

\bdp{5.4.3} \emph{定理.} 对任意 ${}^{L}G$-oper $\mathfrak{F}$，5.1.1 中定义的 $\mathcal{D}$-模 $M_{\mathfrak{F}}$ 有自然的 Hecke $\mathfrak{F}$-特征模结构。

留作习题：函子 $T^i_X$ 与 $H^iV^X\circledast$ 一致（记号见 5.2.4、5.2.5）。于是定理 5.4.3 蕴含 5.2.6。

\bdp{5.4.4} 我们需要 5.4.1–5.4.3 的「带参数」版本。设 $A$ 是交换环。记 $\mathcal{M}(\operatorname{Bun}_G\times X,A)$ 为 $\mathcal{M}(\operatorname{Bun}_G\times X)$ 中 $A$-模构成的范畴（即带 $A$-作用的 $\mathcal{D}$-模）。它有导出版本 $D(\operatorname{Bun}_G\times X,A)$，后者是以 $\mathcal{M}(\operatorname{Bun}_G\times X,A)$ 为心的 t-范畴（见 7.3.13）。如同 5.4.1，范畴 $D(\operatorname{Bun}_G\times X,A)$ 带 $\operatorname{Rep}{}^{L}G$ 的 Hecke 作用。

我们还有 $A\otimes\mathcal{O}_X$-平坦的 $A\otimes\mathcal{D}_X$-模所成张量范畴在 $D(\operatorname{Bun}_G\times X,A)$ 上的显然作用，它与 Hecke 作用交换。因此 $X$ 上带联络的 ${}^{L}G$-丛的任意平坦 $A$-族 $\mathcal{F}_A$ 给出 $\operatorname{Rep}{}^{L}G$ 在 $D(\operatorname{Bun}_G\times X,A)$ 上的作用。

现在对 $M\in\mathcal{M}(\operatorname{Bun}_G,A)$，如同 5.4.2 定义 $M$ 上 Hecke $\mathfrak{F}_A$-特征模结构的概念。下述定理是 5.2.9 的精确版本；由 5.1.2(i)，它蕴含 5.4.3。

\SourcePage{219}

\bdp{5.4.5} \emph{定理.} 5.1.1 中定义的 $\mathcal{D}$-模 $M_{\mathcal{L}}\in\mathcal{M}(\operatorname{Bun}_G,\mathcal{A}_{{}^{L}\mathfrak{g}}(X))$ 有典范的 Hecke $\mathfrak{F}_{\mathcal{L}}$-特征模结构。

\bdp{5.4.6} \emph{注.} (i) 有时（当你想用交换约束时，例如见下一个注或下一节）在超 $\mathcal{D}$-模的框架下处理上述概念是方便的。注意 $\operatorname{Bun}_G$ 上任意 $\mathcal{D}$-模 $M$ 都带典范的 $\mathbb{Z}/2\mathbb{Z}$-分次：$M$ 为偶或奇，取决于 $M$ 支撑在 $\operatorname{Bun}_G$ 的偶分支还是奇分支。记该超 $\mathcal{D}$-模为 $M^{\natural}$。于是 $\natural$ 把 $\mathcal{M}(\operatorname{Bun}_G)$ 与 $\mathcal{M}(\operatorname{Bun}_G)^{\varepsilon}:=\mathcal{M}(\operatorname{Bun}_G)\otimes\operatorname{Vect}^{\varepsilon}$ 的全子范畴 $\mathcal{M}(\operatorname{Bun}_G)^{\natural}$ 等同。对 $D(\operatorname{Bun}_G)$ 与 $D(\operatorname{Bun}_G\times X)$ 同样适用。

$\mathcal{P}$ 在 $D(\operatorname{Bun}_G\times X)$ 上的作用给出 $\mathcal{P}^{\varepsilon}$ 在 $D(\operatorname{Bun}_G\times X)^{\varepsilon}$ 上的作用。$\mathcal{P}^{\natural}\subset\mathcal{P}^{\varepsilon}$ 的作用保持 $D(\operatorname{Bun}_G\times X)^{\natural}$，也保持 $\operatorname{Vect}^{\nabla}(X)$-作用。现在完全如同 5.4.2 定义 $\mathcal{M}(\operatorname{Bun}_G)^{\natural}$ 的 Hecke $\mathfrak{F}$-特征对象的概念。该定义没有带来新东西：$\mathcal{D}$-模 $M$ 是 Hecke $\mathfrak{F}$-特征模，当且仅当 $M^{\natural}$ 是。

(ii) 上述 $\mathcal{M}(\operatorname{Bun}_G)$ 中 $M$ 上的 $\mathcal{F}$-特征模结构定义，用了 $\mathcal{P}$ 上张量结构的卷积构造。可以改用 $\circledast$ 的融合构造重写如下。DOPISAT'!!!
% [Tutor Note: 原书 p.219 此处为俄文注记「DOPISAT'!!!」，意为「待补写」，底本如此，保留原样。]

\bdp{5.4.7} 转向 5.2 中的主要局部定理。我们处于 5.2.12 的设定，故固定 $\mathcal{L}\in Z\operatorname{tors}_{\theta}(\mathcal{O})$，它定义了 $G(K)$ 的中心扩张 $\widetilde{G(K)}=\widetilde{G(K)}_{\mathcal{L}}$，该扩张在群子概形 $G(\mathcal{O})$ 上分裂（见 4.4.9）。我们有相应的扭 Harish-Chandra 模范畴 $\mathcal{M}(\mathfrak{g}\otimes K,G(\mathcal{O}))'$ 与 Harish-Chandra 复形的导出范畴 $D(\mathfrak{g}\otimes K,G(\mathcal{O}))'$（见 7.8.1 与 7.14.1\footnote{故 $1\in\mathbb{C}\subset\mathfrak{g}\otimes K$ 在这些范畴的对象上通过恒等作用。}）。据 7.8.2、7.14.1，$D(\mathfrak{g}\otimes K,G(\mathcal{O}))'$ 带对 $(G(K),G(\mathcal{O}))$ 的 Hecke 幺半范畴 $\mathcal{H}$ 的典范作用 $\circledast$。由于 $\mathcal{P}$ 是 $\mathcal{H}$（的心）的幺半子范畴，我们的 $D(\mathfrak{g}\otimes K,G(\mathcal{O}))'$ 是 $\mathcal{P}$-模。
% 原书 p.219 末句完结，但 5.4.7 的论述续于 p.220（下一单元）。
